\chapter{Exportación}
\label{sec:man-export}

\section{IOCs}

\emph{Archivo \textrightarrow{} Exportar \textrightarrow{} IOCs}.
Parámetros:

\begin{itemize}
    \item \textbf{Filtro por tag} (ej.\ \code{ioc}): exporta sólo
    entidades con ese tag. Default: todas las IPs, domains, hashes
    marcados en la sesión.
    \item \textbf{Formato}: STIX 2.1 (default), JSON plano, CSV.
\end{itemize}

El output lista las entidades como indicadores STIX con tipo
(\code{ipv4-addr}, \code{domain-name}, \code{file:hashes.SHA256}),
primer y último visto, y relaciones explícitas si las agregaste
vía tags.

\section{Publicar IOCs a Sentinel}

Si tu organización usa Microsoft Sentinel:
\emph{Integraciones \textrightarrow{} Publicar IOCs}. Requiere:

\begin{itemize}
    \item \code{GRAPHHUNTER_IOC_PUBLISH_ENABLED=1} configurada.
    \item Credenciales Azure configuradas (ver
    cap.~\ref{sec:man-integraciones}).
\end{itemize}

La app pide confirmación antes de publicar; podés elegir
\emph{dry-run} para ver qué se enviaría sin enviar.

\section{OCSF}

Exporta los eventos canonical ya normalizados (ver \emph{Canonical
OCSF} en Especificación):

\begin{itemize}
    \item \textbf{Scope}: toda la sesión, o un subset filtrado por
    ventana temporal, tipos de entidad, o tags.
    \item \textbf{Formato}: JSONL (una línea por evento) o JSON
    array.
\end{itemize}

Útil para:

\begin{itemize}
    \item Reingestar en otra plataforma que hable OCSF.
    \item Archivo a largo plazo en formato estandarizado.
    \item Compartir con equipos que usan otras herramientas.
\end{itemize}

\section{Subgrafos}

\emph{Archivo \textrightarrow{} Exportar \textrightarrow{}
Subgraph} a partir de la selección actual (del modo explorador) o
de un path específico (del modo caza). Formatos:

\begin{itemize}
    \item \textbf{GraphSON 3.0}: compatible con TinkerPop /
    Gremlin.
    \item \textbf{JSONL}: lista de vértices y aristas, simple.
    \item \textbf{DOT}: para renderizar con Graphviz.
\end{itemize}

\section{Resultados de hunt}

Un hunt con sus paths y scores se exporta con
\emph{Hunt \textrightarrow{} Export results}. Formatos CSV o JSON.
CSV incluye una columna por componente de score más el compuesto,
útil para análisis en planilla.

\section{Sigma}

Una hipótesis útil puede guardarse como regla Sigma
(\emph{Hunt \textrightarrow{} Save as Sigma}). La conversión mapea
los nodos y relaciones del DSL a la estructura \code{detection} de
Sigma. No todo DSL es convertible 1:1 — hipótesis que dependen de
ventanas temporales o predicados exóticos emiten Sigma con notas
sobre qué parte no se tradujo.

\section{Abrir ticket}

\emph{Integraciones \textrightarrow{} Crear ticket}. Requiere
\code{GRAPHHUNTER_TICKET_PUBLISH_ENABLED=1} y credenciales de Jira
o ServiceNow. El ticket incluye:

\begin{itemize}
    \item Título y descripción editables.
    \item Paths seleccionados con su desglose de scoring.
    \item IOCs tagueados.
    \item Link a la sesión si está publicada.
\end{itemize}

Igual que IOCs, soporta \emph{dry-run}.

\section{Workflow típico de entrega}

\begin{enumerate}
    \item Taguear entidades y paths relevantes durante el
    análisis.
    \item \emph{Exportar IOCs} filtrado por tag — lo mandas al SOC
    o lo publicas a Sentinel.
    \item \emph{Exportar subgrafo} de los paths clave — útil para
    ilustrar en el reporte.
    \item \emph{Crear ticket} con narrativa completa y paths
    anexados.
    \item Mantener la sesión archivada por si se necesita volver.
\end{enumerate}
